% 第一章第二节：沿事务步骤解释余额检查、提交与回滚。
四张表的关系已经明确。现在沿着 Alice 的购买检查具体写入：从 10000 分扣除 3000 分，得到购买 P1、订单 O1、流水 D1 和 7000 分余额。先逐步看 SQL，再讨论应用怎样组织事务和返回结果。

下面使用固定 UUID 演算，方便逐句辨认记录之间的关联。每组演算都说明自己的初始状态。后续从 HTTP 入口重新发起购买时，服务会生成新的购买编号，数据也会随这次新购买继续变化。

\subsection*{先确定账户与价格，再开始写入}

客户端发送的请求只有商品编号：

\begin{BookCode}
{"product_id": "paper-001"}
\end{BookCode}

服务收到商品编号后，先检查目录里是否有这件商品，再从目录取得价格 3000 和币种 \code{CNY}。商品目录可以保存在数据库中，也可以由其他受控配置提供。本章只需要两条目录数据：\code{paper-001} 售价 3000 分，\code{paper-002} 售价 8000 分。

账户来自已经验证的调用方身份。服务先验证登录凭据，确认用户，再取得该用户有权使用的账户。Alice 的身份对应账户 1，因此传入购买事务的也是账户 1。由服务完成这一步映射，账户权限就能在业务写入之前得到检查。

价格由服务端目录确定。客户端选择商品，服务端据此取得应付金额；如果直接采用客户端填写的金额，调用方就能自行改价。因此，这个购买接口的请求体只有 \code{product\_id}，服务按约定检查字段，格式错误或带有额外字段的请求返回 400。

这些检查完成后，应用准备好一份输入：账户 1、商品 \code{paper-001}、金额 3000、币种 CNY，以及新生成的购买、订单和扣款 UUID。整个事务重复使用这份输入。扣款后不再回头读取一次目录价格，防止订单金额与刚才实际扣除的金额不同。

\subsection*{用一笔事务组织整组变化}

如果把四张表各自提交，会出现难以解释的中间结果。先扣款并提交，再建订单；一旦建订单失败，账户已经少了钱。先建已付款订单并提交，再扣款；一旦余额不足，数据库却已有一张已付款订单。

因此，这几项写入要么一起提交，要么一起撤回。在 PostgreSQL 中，用 \code{BEGIN} 开始事务，\code{COMMIT} 提交，\code{ROLLBACK} 撤回。连接在事务内部能够读取自己的写入，其他连接的普通查询读不到这些尚未提交的改动。

先记下购买的输入，表示这笔事务正在处理 P1：

\begin{BookCode}
BEGIN;

INSERT INTO purchases (
    purchase_id, account_id, product_id,
    amount_minor, currency, status
) VALUES (
    '10000000-0000-0000-0000-000000000001',
    1, 'paper-001', 3000, 'CNY', 'pending'
);
\end{BookCode}

这里给出一个完整 UUID，后面用同一 UUID 表示 P1。订单 O1 和流水 D1 分别使用以 2000、3000 开头的示例 UUID，便于在 SQL 中辨认。

此时事务内部已有 P1，账户仍为 10000，订单和流水尚未建立。发起写入的连接可以查到自己刚插入的 P1；另一个连接的普通查询仍只能读到已经提交的数据。图~\ref{fig:ch01-visibility} 把这两个连接放在同一条时间线上：事务内部正在形成购买，外部查询要等到提交以后才能看到它。

\begin{figure}[htbp]
\centering
\begin{tikzpicture}[x=1cm,y=1cm]
\node[booknode,minimum width=4cm] at (2.2,0) {连接 A：正在处理购买};
\node[booknode,minimum width=4cm] at (10.3,0) {连接 B：普通查询};
\draw[draw=Muted!60,dashed] (2.2,-.4)--(2.2,-5.6);
\draw[draw=Muted!60,dashed] (10.3,-.4)--(10.3,-5.6);
\node[booknode,text width=4.2cm] (insert) at (2.2,-1.4)
  {BEGIN，插入 P1\\A 自己可以查到 P1};
\node[booknode,text width=4.2cm] (before) at (10.3,-1.4)
  {查询 P1\\结果为空：P1 尚未提交};
\draw[draw=Muted!60,dotted] (insert.east)--(before.west);
\node[booknode,text width=4.2cm] (writes) at (2.2,-3)
  {继续扣款、建立订单与流水\\更新购买结果};
\draw[bookarrow] (insert)--(writes);
\draw[bookalert] (0,-4.2)--(12.5,-4.2);
\node[fill=white,font=\footnotesize\sffamily,text=Accent] at (6.3,-4.2)
  {A 提交完成};
\node[bookresult,text width=4.2cm] at (10.3,-5.1)
  {提交后重新查询\\读到 P1、订单与流水};
\end{tikzpicture}
\caption{同一笔购买在提交前后的可见性。图中连接 B 使用 Read Committed 的普通查询。}
\label{fig:ch01-visibility}
\end{figure}


\subsection*{在同一条语句里检查余额并扣款}

扣款需要账户存在、余额足够。把余额检查直接写进更新语句：

\begin{BookCode}
UPDATE accounts
SET balance_minor = balance_minor - 3000
WHERE account_id = 1
  AND balance_minor >= 3000
RETURNING balance_minor;
\end{BookCode}

Alice 的余额原为 10000，满足条件。数据库执行 \code{10000 - 3000}，将事务内部的余额更新成 7000，并返回一行 7000。\code{RETURNING} 让应用直接拿到这次更新产生的新值，无须再发一条查询去猜刚才扣款后的余额。

如果余额原为 2000，条件就不成立，语句更新零行，返回结果也为空。应用据此进入余额不足分支。这个判断依赖账户存在这一前提：本章账户预先建立，运行期间保持存在；若应用传入了不存在的账户，前面的购买插入就会因外键约束失败，进入错误处理。

也可以设想另一种写法：先查询余额，由应用判断金额是否足够，再把算好的余额写回。两个购买都读到余额 10000，各自要扣 7000，就可能都通过检查，并各自算出 3000。如果两边随后都无条件写入 3000，两个订单可能都成功，账户却只减少了 7000。非负约束会接受这个数值，因为 3000 仍然大于零。

条件更新把余额检查和余额扣减放在同一次行更新中。同一账户上的另一次更新可能等待当前事务结束。在 \code{READ COMMITTED} 隔离级别下，如果前一个更新已经提交，等待的事务会在更新后的行版本上重新判断 \code{WHERE} 条件。前一笔扣掉 7000 后只剩 3000，后一笔再扣 7000 就无法满足条件。

图~\ref{fig:ch01-concurrent} 画出了先后两笔更新。这里使用的是 PostgreSQL 对更新冲突的处理规则。在 \code{READ COMMITTED} 下，普通查询则由每条语句取得自己的读快照，两条查询之间可能出现新的提交。稍后核对余额与流水时，我们会在同一条查询里读取它们，使两边采用同一个语句快照。

\begin{figure}[htbp]
\centering
\begin{tikzpicture}[x=1cm,y=1cm]
\node[booknote] at (6,1.45) {初始余额 10000 分；两笔不同购买各需 7000 分};
\node[booknode,text width=4cm] (a1) at (2.5,0) {购买 A 先更新账户\\条件成立，余额改为 3000};
\node[booknode,text width=4cm] (b1) at (9.5,0) {购买 B 也请求更新\\等待账户行上的更新冲突解除};
\node[bookresult,text width=4cm] (a2) at (2.5,-2) {A 补齐订单与流水\\提交，余额 3000 分生效};
\node[booknode,text width=4cm] (b2) at (9.5,-2) {A 提交后，B 继续\\按新行版本重查条件};
\node[bookresult,text width=4cm] (b3) at (9.5,-4) {3000 小于 7000\\更新零行\\B 保存余额不足};
\draw[bookarrow] (a1)--(a2);
\draw[bookarrow] (b1)--(b2);
\draw[bookalert] (a2.east)--node[above,font=\footnotesize]{提交后继续}(b2.west);
\draw[bookarrow] (b2)--(b3);
\node[booknote,text width=4cm] at (2.5,-4) {最终：仅 A 扣款 7000 分\\账户余额 3000 分};
\end{tikzpicture}
\caption{在 PostgreSQL 的 Read Committed 隔离级别下，等待更新的事务会基于已经提交的新行版本重新判断条件。}
\label{fig:ch01-concurrent}
\end{figure}


\subsection*{扣款成功后补齐订单和流水}

扣款语句返回 7000 后，还不能给 Alice 发成功响应。现在只在当前事务内部改了余额，订单和流水都没有建立。先把成交内容写进订单：

\begin{BookCode}
INSERT INTO orders (
    order_id, purchase_id, product_id,
    amount_minor, currency, status
) VALUES (
    '20000000-0000-0000-0000-000000000001',
    '10000000-0000-0000-0000-000000000001',
    'paper-001', 3000, 'CNY', 'paid'
);
\end{BookCode}

这条订单把 O1 指向 P1，商品与金额沿用刚才那份输入。此时事务内部有账户 7000、购买 P1 和订单 O1，仍然缺少扣款依据。接着写入流水：

\begin{BookCode}
INSERT INTO debits (
    debit_id, order_id, account_id,
    amount_minor, balance_after_minor
) VALUES (
    '30000000-0000-0000-0000-000000000001',
    '20000000-0000-0000-0000-000000000001',
    1, 3000, 7000
);
\end{BookCode}

D1 引用 O1，账户为 1，扣款额为 3000，扣后余额是刚才 \code{RETURNING} 得到的 7000。流水插入成功后，购买所需的业务记录齐了，把 P1 改为成功并提交：

\begin{BookCode}
UPDATE purchases
SET status = 'succeeded'
WHERE purchase_id = '10000000-0000-0000-0000-000000000001';

COMMIT;
\end{BookCode}

提交以后，其他连接可以查到这一整组已提交数据。余额减少、订单建立和流水写入随同一笔事务生效。它们在执行过程中出现的中间状态，仅供当前事务内部使用。

四个地方的金额值得再对照一次：账户减少 3000，P1 记录金额 3000，O1 记录成交额 3000，D1 记录扣款额 3000。这个对应由相同输入和同一个事务共同维护。数值相等让账对得上，事务让它们一起保存下来。

\subsection*{余额不足也走到一个确定结果}

Alice 现在剩下 7000 分。她接着购买 \code{paper-002}，服务端价格为 8000 分。服务为这次受理生成新的购买编号 P2，再开始一笔新事务。先插入 P2，金额 8000，状态 \code{pending}，然后执行：

\begin{BookCode}
UPDATE accounts
SET balance_minor = balance_minor - 8000
WHERE account_id = 1
  AND balance_minor >= 8000
RETURNING balance_minor;
\end{BookCode}

7000 小于 8000，数据库没有修改账户，返回零行。应用跳过订单和流水插入，把 P2 的状态改为 \code{insufficient\_funds} 后提交。账户仍为 7000，P1、O1、D1 不变，新增加的只有 P2。

图~\ref{fig:ch01-branches} 将余额不足与成功分支并列放在一起：两个分支都保存受理结果，订单与流水只出现在扣款成功的一边。

余额不足是正常的业务分支。数据库能够执行语句，服务也知道为什么没有扣款，因此可以保存这个结果，再把它返回给调用方。数据库连接失败、SQL 拼错或存储错误属于另一类处理错误，不能一律转成 \code{insufficient\_funds}。否则，调用方会把系统暂时无法处理误解为自己的钱不够。

P2 提交后就以余额不足的结果结束。Alice 以后增加账户余额，再次点击购买，服务会受理另一笔购买，并按当时的余额重新判断。P2 则继续保留原来的结果，供详情查询使用。

\begin{figure}[htbp]
\centering
\begin{tikzpicture}[x=1cm,y=1cm]
\node[booknode,text width=5.4cm] (p) at (6,0) {开始事务，插入购买记录\\状态为 pending};
\node[booknode,text width=5.4cm] (u) at (6,-1.5) {执行带余额条件的 UPDATE\\检查返回结果};
\node[booknode,text width=4.4cm] (yes) at (2.3,-3.5) {返回一行：扣款成功\\建立订单与流水\\购买改为 succeeded};
\node[booknode,text width=4.4cm] (no) at (9.7,-3.5) {返回零行：余额不足\\账户、订单与流水保持原样\\购买改为 insufficient\_funds};
\node[bookresult,text width=5.4cm] (c) at (6,-5.6) {提交这一笔事务\\保存对应分支的完整结果};
\draw[bookarrow] (p)--(u);
\draw[bookarrow] (u.south)--++(0,-.35)-|node[above,pos=.75,font=\footnotesize]{一行}(yes.north);
\draw[bookarrow] (u.south)--++(0,-.35)-|node[above,pos=.75,font=\footnotesize]{零行}(no.north);
\draw[bookarrow] (yes.south)|-(c.west);
\draw[bookarrow] (no.south)|-(c.east);
\end{tikzpicture}
\caption{扣款成功与余额不足分别形成确定的业务结果。执行错误另走回滚处理。}
\label{fig:ch01-branches}
\end{figure}


\subsection*{写到一半出错时，整笔事务回滚}

再从一个全新的 10000 分账户出发，检查成功路径的中途失败。手工执行时可以把最后的 \code{COMMIT} 换成 \code{ROLLBACK}，观察已经执行的步骤怎样被撤回。

第一处停在购买插入之后。事务内部只有 \code{pending} 购买，还没有扣款。回滚后，购买记录消失，余额保持 10000，订单和流水仍为空。

第二处停在条件扣款之后。事务内部已有购买，余额显示 7000，还没有订单。回滚会同时撤回购买插入和余额更新。重新用一个正常查询读取，应当回到初始 10000 分。

第三处停在订单插入之后、流水插入之前。事务内部能看到购买、7000 分余额和订单。回滚会一并撤回订单插入、余额更新和购买插入，账户余额恢复为 10000。

如果流水已经插入，随后更新购买结果的语句失败，整笔事务仍应回滚。购买成功的条件包含订单、扣款与购买结果共同提交；忽略其中一次写入错误，会让返回结果与实际记录发生偏离。

图~\ref{fig:ch01-rollback} 汇集了这几处中断。若应用自行执行反向操作，例如在失败后给账户加回 3000，就需要继续处理反向更新的失败、并发修改和重复执行。对于这组尚未提交的本地写入，数据库回滚已经能够撤销此前的变化。应用应当把它们放在同一事务中，由数据库完成撤回。

\begin{figure}[htbp]
\centering
\begin{tikzpicture}[x=1cm,y=1cm]
\node[booknote] at (5.8,.65) {同一笔事务提交之前};
\node[booknode,text width=5cm] (s1) at (2.6,-.3) {插入购买后\\余额 10000，购买 pending};
\node[booknode,text width=5cm] (s2) at (2.6,-1.9) {条件扣款后\\余额 7000，购买 pending};
\node[booknode,text width=5cm] (s3) at (2.6,-3.5) {插入订单后\\余额 7000，购买与订单已写入};
\node[booknode,text width=5cm] (s4) at (2.6,-5.1) {插入流水后\\尚待更新购买结果并提交};
\node[bookresult,text width=4.3cm,minimum height=3cm] (r) at (10,-2.7)
  {任何一处中止并回滚\\[7pt]余额恢复为 10000\\本次购买、订单、流水\\全部撤回};
\draw[bookarrow] (s1)--(s2);
\draw[bookarrow] (s2)--(s3);
\draw[bookarrow] (s3)--(s4);
\foreach \s in {s1,s2,s3,s4}
  \draw[bookalert] (\s.east)--(r.west);
\end{tikzpicture}
\caption{未提交的中间状态可以各不相同，回滚会撤销本次事务已经执行的整组写入。}
\label{fig:ch01-rollback}
\end{figure}


\subsection*{把事务提交放在响应之前}

现在把前面的 SQL 连成应用的控制流程。进入事务之前，先确定已经授权的账户、商品、成交金额和币种，并生成本次购买使用的编号。下面的伪代码只表示步骤与先后关系：

\begin{BookCode}
输入 <- 已授权的账户、商品、成交金额、币种
编号 <- 生成购买、订单和扣款编号
开始事务
    写入购买记录，状态为 pending
    扣款结果 <- 按余额条件更新账户，取得扣后余额
    如果扣款结果含有一行
        写入已付款订单
        写入扣款流水，保存扣后余额
        将购买结果记为 succeeded
    否则
        将购买结果记为 insufficient_funds
    组装待返回的购买结果
提交事务
确认提交完成
返回购买结果
\end{BookCode}

条件更新返回的一行数据带有扣后余额，成功分支把这个值写入流水。返回零行时，购买以余额不足结束。两条分支都保存购买结果，订单和流水则只在扣款成功的分支中建立。

所有数据库操作使用同一连接，并处在同一笔显式事务中。实现时应核对数据库接口的事务行为：若每条语句执行后就单独提交，前面讨论的整体回滚就失去了前提。SQL 排在同一个函数里，并不能自动形成事务。

事务采用前面说明的 \code{READ COMMITTED} 隔离级别。应用执行带条件的余额更新，数据库负责处理同一账户上的更新冲突。成交金额和各项编号在事务开始前已经确定，后续写入一直使用这份输入。

组装购买结果时，应用可以先在内存中准备返回字段。随后提交事务，等待数据库确认，再把结果交给 HTTP 层发送。这样，调用方收到成功响应时，对应的购买、订单和扣款已经保存。

事务内部某条 SQL 执行失败时，应中止当前处理并回滚。提交阶段如果连接中断，应用还需要核实数据库是否已经提交；单凭连接错误无法确定结果。下一章将从提交与响应之间的这个间隔继续讨论。

\subsection*{对照业务记录，再核对余额}

处理完成后，先沿购买查订单和流水，检查三个编号是否连得上、商品与金额是否一致。再从账户的整体变化核对流水合计。对于本章初始 10000 分、只有购买扣款的账户，有：

\[
  \mbox{当前余额}
  = 10000 - \mbox{已提交扣款流水金额之和}
\]

下面在一条查询中读取余额与扣款合计：

\begin{BookCode}
SELECT a.account_id, a.balance_minor,
       COALESCE(SUM(d.amount_minor), 0) AS debited_minor,
       10000 - COALESCE(SUM(d.amount_minor), 0)
           AS expected_balance_minor
FROM accounts AS a
LEFT JOIN debits AS d ON d.account_id = a.account_id
WHERE a.account_id = 1
GROUP BY a.account_id, a.balance_minor;
\end{BookCode}

\code{LEFT JOIN} 保留还没有扣款的账户，\code{COALESCE} 把空合计转换成零。购买前得到当前余额 10000、扣款合计零、预期余额 10000；一次 3000 分购买成功后，得到 7000、3000、7000。随后那次 8000 分购买余额不足，三个值保持不变。

这项核对覆盖账户的全部扣款历史。以最初的 10000 分为起点，就需要累计自那时起的每一条扣款流水。后续若增加充值、退款和人工调整，也应把它们对余额的影响计入核对公式。
